\documentclass[11pt,a4paper]{article}

% Packages
\usepackage[utf8]{inputenc}
\usepackage[T1]{fontenc}
\usepackage[margin=0.75in]{geometry}
\usepackage{enumitem}
\usepackage{hyperref}
\usepackage{titlesec}
\usepackage{xcolor}

% Formatting
\pagestyle{empty}
\setlength{\parindent}{0pt}
\setlist[itemize]{leftmargin=*, noitemsep, topsep=3pt}

% Section formatting
\titleformat{\section}{\Large\bfseries}{}{0em}{}[\titlerule]
\titlespacing*{\section}{0pt}{12pt}{6pt}

\titleformat{\subsection}[runin]{\bfseries}{}{0em}{}
\titlespacing*{\subsection}{0pt}{8pt}{0pt}

% Colors
\definecolor{linkcolor}{RGB}{0,102,204}
\hypersetup{
    colorlinks=true,
    urlcolor=linkcolor
}

\begin{document}

% Header
\begin{center}
    {\LARGE \textbf{ERIC MAYEUX}}\\
    \vspace{4pt}
    {\large Chef de projet Implémentation logiciels | Software Implementation Project Manager}\\
    \vspace{4pt}
    Montréal, QC, Canada\\
    \href{mailto:mayeux.e@gmail.com}{mayeux.e@gmail.com} | 438-884-7645
\end{center}

\vspace{8pt}

\section*{Profil Professionnel}
\noindent
\textbf{Chef de projet en implémentation de logiciels} avec 10+ ans d'expérience en \textbf{gestion de projets d'implémentation} TI dans le secteur bancaire et financier. Expert en \textbf{gestion simultanée de plusieurs projets}, coordination d'équipes multidisciplinaires et \textbf{gestion des parties prenantes} internes et externes. Maîtrise de la \textbf{méthodologie Agile} et des cadres hybrides. Certifié \textbf{PMP} (en cours) et CSM. Bilinguisme français/anglais.

\section*{Compétences Techniques}

\textbf{Gestion de projets d'implémentation:} Planification, jalons, livrables, calendriers fiables, gestion simultanée de plusieurs projets, budgets (2M\$+)

\textbf{Coordination \& Parties prenantes:} Équipes multidisciplinaires, consultants externes, fournisseurs, communication efficace, gestion de la relation client

\textbf{Gestion des risques:} Anticipation, plans de contingence, actions correctives, suivi hebdomadaire d'avancement

\textbf{Outils:} Jira, MS Project, Confluence, Datadog, Splunk, tableaux de bord KPI

\textbf{Méthodologies:} Agile, Scrum, SAFe, Kanban, Waterfall, hybride, PMP

\textbf{Intelligence Artificielle:} Agents IA, automatisation, Playwright MCP, Crew AI

\textbf{Architecture \& Cloud:} AWS Cloud Practitioner, Architectures Distribuées, Big Data

\vspace{6pt}

\section*{Réalisations}
\begin{itemize}
    \item \textbf{Gestion simultanée de plusieurs projets} d'implémentation TI à la BNC (secteur bancaire)
    \item Pilotage des releases : de l'estimation à la \textbf{livraison} en respectant délais et budgets
    \item Coordination d'\textbf{équipes multidisciplinaires} (QA, Dev, PO, Analystes, équipes Thaïlande)
    \item \textbf{Rapports hebdomadaires d'avancement} aux parties prenantes internes et haute direction
    \item Mise en place de cadres de \textbf{gestion des risques} et métriques qualité (Jira, Datadog)
    \item \textbf{Coaching} et orientation des membres moins expérimentés ; recrutement et onboarding
\end{itemize}

\section*{Expérience Professionnelle}

\subsection*{Scrum Master -- Gestion de projet | Project Manager} \hfill \textit{Novembre 2025 -- Présent}\\
\textbf{Banque Nationale du Canada (BNC)} -- Montréal, QC\\
Pilote les projets d'implémentation TI dans un contexte bancaire agile (SAFe).
\begin{itemize}
    \item Gère la \textbf{planification, coordination et livraison} des projets TI de A à Z
    \item Coordonne des \textbf{équipes multidisciplinaires} ; assure l'alignement sur les livrables
    \item Anticipe les \textbf{risques} et plans de contingence ; produit des \textbf{rapports d'avancement} hebdomadaires
    \item Maintient les \textbf{parties prenantes} informées et aligne les équipes techniques et fonctionnelles
    \item \textbf{Oriente et coache} les membres de l'équipe sur les pratiques de gestion de projet
    \item Développe des agents IA pour l'\textbf{automatisation} des tâches de gestion
\end{itemize}

\subsection*{Intégrateur Sénior Qualité | Senior Quality Integrator} \hfill \textit{Octobre 2023 -- Novembre 2025}\\
\textbf{Banque Nationale du Canada (BNC)} -- Montréal, QC
\begin{itemize}
    \item \textbf{Gestion de plusieurs projets} simultanés d'implémentation et d'amélioration qualité
    \item Coordination multi-sites (Thaïlande) : planification, livrables, \textbf{gestion des risques}
    \item \textbf{Coaching} et mentorat des QA ; recrutement et formation continue
\end{itemize}

\subsection*{Intégrateur Qualité -- SDET | Quality Integrator -- Software Development Engineer in Test} \hfill \textit{Janvier 2022 -- Octobre 2023}\\
\textbf{Banque Nationale du Canada (BNC)} via Groupe Astek -- Montréal, QC
\begin{itemize}
    \item Conception et suivi de la \textbf{stratégie de test} pour des projets d'implémentation logicielle
    \item Stack : Selenium, AppliTools, RestAssured, JMeter, Gatling, K6
\end{itemize}

\subsection*{QA Automaticien | QA Automation Engineer} \hfill \textit{Juin 2021 -- Décembre 2021}\\
\textbf{AODocs} -- Paris, France
\begin{itemize}
    \item \textbf{Implémentation} des tests non-régression ; POC Flutter ; rapports Allure
\end{itemize}

\subsection*{Automaticien CI/CD | DevOps Automation Engineer} \hfill \textit{Janvier 2020 -- Juin 2021}\\
\textbf{ENGIE DIGITAL} via Groupe Hénix -- Paris, France
\begin{itemize}
    \item Création CI/CD et \textbf{implémentation des outils} : Jenkins, Docker ; pilotage Jira, KPIs
\end{itemize}

\subsection*{Product Owner} \hfill \textit{Juin 2019 -- Septembre 2019}\\
\textbf{Hénix} -- Montrouge, France
\begin{itemize}
    \item \textbf{Implémentation} de plugins Jira/Squash ; gestion du backlog et priorisation
\end{itemize}

\subsection*{Développeur Logiciel | Software Developer} \hfill \textit{Mai 2017 -- Mai 2019}\\
\textbf{MEDIAPOST} via Groupe Hénix -- Paris, France
\begin{itemize}
    \item Développement back office, webservices SOAP/REST, scripts BDD ; création d'application interne
\end{itemize}

\subsection*{Consultant QA \& Automatisation | QA Automation Consultant} \hfill \textit{Avril 2016 -- Avril 2017}\\
\textbf{ENGIE IT} via Groupe Hénix -- Saint-Ouen, France
\begin{itemize}
    \item Pilotage opérationnel, coordination fournisseurs, administration HP ALM/UFT
\end{itemize}

\section*{Certifications \& Formations}

\textbf{PMP Certification} (en cours) \hfill \textit{2026}

Project Management Professional -- Project Management Institute (PMI)

\textbf{AWS Cloud Practitioner} \hfill \textit{2026}

Amazon Web Services (AWS)

\textbf{Certified Scrum Master (CSM)} \hfill \textit{2024}

Scrum Alliance

\textbf{Jira ACP-600 -- Project Administration in Jira Server} \hfill \textit{2019}

Atlassian Certified Professional

\textbf{Cursus Automaticien et DevOps} \hfill \textit{2019}\\
\textit{AFCEPF -- Montrouge, France}

\textbf{Cursus Architecture logiciel \& Analyste informaticien - Certifié Master II} \hfill \textit{2017}\\
\textit{AFCEPF-HENIX -- Malakoff, France}

\section*{Formation Académique}

\textbf{Licence en Gestion (Bachelor's Degree in Management)} \hfill \textit{2011}\\
École Supérieure de Gestion (E.S.G) -- Paris, France

\section*{Compétences Linguistiques}

\textbf{Français:} Langue maternelle (Native)\\
\textbf{Anglais:} Niveau Avancé (Advanced / Professional Working Proficiency)

\end{document}
